\begin{proofbox}
	\textbf{\textcolor{CProof}{思路提示}}

	采用反证法。假设 $\alpha$ 与 $\beta$ 不平行，则它们相交于一条直线。利用线面平行性质，将 $a,b$ 平行于 $\beta$ 转化为 $a,b$ 平行于交线，从而导出 $a\parallel b$，与已知 $a,b$ 相交矛盾，故原结论成立。

	\medskip
	\textbf{\textcolor{CProof}{正式推导}}
	\begin{description}[style=nextline, leftmargin=2.8em, labelsep=0.5em, itemsep=0.55em]
		\item[\textbf{步骤一（假设相交）：}] 假设平面 $\alpha$ 与 $\beta$ 不平行，则它们必有一条交线，设 $l=\alpha\cap\beta$。
			\[
				\alpha\cap\beta = l
			\]
			\par\textbf{理由：}
			两平面若不平行则相交于一条直线（立体几何公理）。

		\item[\textbf{步骤二（$a\parallel l$）：}] 由 $a\parallel\beta$ 且 $a\subset\alpha$，结合线面平行性质定理，得 $a\parallel l$。
			\[
				a\parallel\beta,\; a\subset\alpha,\; l=\alpha\cap\beta \;\Longrightarrow\; a\parallel l
			\]
			\par\textbf{理由：}
			若一条直线平行于一个平面，且过该直线的平面与这个平面相交，则该直线与交线平行。

		\item[\textbf{步骤三（$b\parallel l$）：}] 同理，由 $b\parallel\beta$ 且 $b\subset\alpha$，可得 $b\parallel l$。
			\[
				b\parallel\beta,\; b\subset\alpha,\; l=\alpha\cap\beta \;\Longrightarrow\; b\parallel l
			\]
			\par\textbf{理由：}
			与步骤二相同，应用线面平行性质定理。

		\item[\textbf{步骤四（导出矛盾）：}] 由 $a\parallel l$ 且 $b\parallel l$ 可得 $a\parallel b$，这与已知 $a\cap b=P$ 矛盾。
			\[
				\left.\begin{array}{l}
				a\parallel l\\
				b\parallel l
		\end{array}\right\} \;\Longrightarrow\; a\parallel b
		\]
		\par\textbf{理由：}
		平行于同一直线的两直线平行（空间成立）。因此 $a\parallel b$，但已知 $a\cap b=P$，矛盾，故假设不成立。
\end{description}

\medskip
\textbf{\textcolor{CProof}{结论回扣}}

\textbf{由反证法，假设“$\alpha$ 与 $\beta$ 不平行”导致矛盾，故 $\alpha\parallel\beta$。}
\end{proofbox}
